% presentation
\documentclass{beamer}

\usetheme{Boadilla}

% rus lang
\usepackage[main=russian,english]{babel}

% insert images
\usepackage{wrapfig}

\title{Лекция 2. Линейная регрессия}
\subtitle{Основы интеллектуального анализа данных}
\author{Полузёров Т. Д.}
\institute{БГУ ФПМИ}
\date{}

\begin{document}
	
	\begin{frame}
		\titlepage
	\end{frame}

		
	\begin{frame}
		\frametitle{Опрделение линейной регрессии}
		
		Модель линейной регрессии
		
		$$
		a(x) = \omega_0 + \sum_{j=1}^{n} \omega_j x_j = 
		\omega_0 + \langle \omega, x\rangle 
		$$
		Определяется вектором коэффициентов $\omega = (\omega_1, ..., \omega_n) \in \mathbb{R}^{n}$ и свободным членом $\omega_0 \in \mathbb{R}$
		
		\vspace{15}
		
		Для упрощения формул добавим к признаковому описанию объектов признак равный единице $x := (1, x_1, ..., x_n), \omega=(\omega_0, \omega_1, ..., \omega_{n})$
		
		$$
		a(x) = \langle \omega, x \rangle
		$$
		
	\end{frame}


	\begin{frame}
		\frametitle{Метод наименьших квадратов (МНК)}
		Пусть задан датасет $(X, y)_{i=1}^{\ell}$, где $X = (x_i)_{i=1}^{\ell} \subseteq \mathbb{R}^{\ell \times n}$ - матрица признаков, $y = (y_i)_{i=1}^{\ell}\subseteq \mathbb{R}^{\ell}$ - вектор значений
		
		Мы хотим приблизить эту зависимость моделью вида: $a(x) = \langle \omega, x \rangle$
		
		Ошибка на объекте $x_i$:
		
		$$
		e_i = a(x_i) - y_i
		$$
		
		Подберем коэффициенты $\omega$ таким образом:
		$$
		\mathcal{L}(a, x, y) = (a(x_i) - y_i)^{2} = e_i^{2}
		$$ - квадратичная функция потерь
		
		$$
		Q(a, X, y) 
		= \frac{1}{\ell} \sum_{i=1}^{\ell} \mathcal{L}(a, x_i, y_i)
		= \frac{1}{\ell} \sum_{i=1}^{\ell} (a(x_i) - y_i)^{2}
		= \frac{1}{\ell} \sum_{i=1}^{\ell} e_i
		\to \min_{\omega}
		$$
	\end{frame}

	\begin{frame}
		\frametitle{Решение МНК}
		Точное решение:
		$$
		\omega = (X^T X)^{-1} X^Ty
		$$
		
		
		
		
	\end{frame}




\end{document}